\clearpage
\section{Actual composition of translated contexts}
\label{sec:copying-translated-composition}
\termdef{copying-translated-composition}
The constructor interpretation already respects composition of annotated term
contexts. We now connect this identity to actual replacement at certified graph
holes. Two independently translated contexts, with the inner host padded to
fit its outer occurrence, compose to a graph recursively isomorphic to the
direct translation of the composed context. Filling its selected hole with
any suitably padded graph also gives the direct constructor interpretation.
The isomorphisms have equal outer depths; no additional common-allowance
comparison or execution hypothesis is used.

\subsection{Certified open graph contexts compose}
Let $A$ have a sole hole accepting the outer interface of $B$, and let $B$
have a sole hole accepting the argument graph. Form $A\circ B$ by replacing
the hole of $A$ with the graph of $B$. Its selected occurrence is the address
of the original hole of $B$, followed through that replacement. The existing
sole-hole composition theorem proves that it is still the only hole.
Both contexts have binding-free certificates. Since the hole input is the
singleton term interface, it is sourced; the proved binding-free plugging
closure therefore certifies the composite. This use respects the zero-input
qualification of Section~\ref{sec:copying-binding}.
\lean{Copying/TranslatedComposition}{10}{17}{GraphContext.compose}

For every graph $G$ at the local allowance, addressed plugging associativity
gives a recursive isomorphism
\[
  (A\circ B)[G]\ \cong\ A[B[G]].
\]
It retains all nested boxes and ordered interfaces, allowing the carrier
reassociation already handled by the generic graph proof. Specialising $G$
to the graph of an open program yields the corresponding equality up to
recursive isomorphism between composed filling and successive filling.
The open-program certificates do not add operational premises.
\lean{Copying/TranslatedComposition}{20}{24}{GraphContext.compose_plug_iso}
\lean{Copying/TranslatedComposition}{26}{30}{GraphContext.compose_fill_iso}

The exact depth function also composes:
$(C\circ D).\operatorname{depthAt}(d)
=C.\operatorname{depthAt}(D.\operatorname{depthAt}(d))$.
Induction on the outer context proves this identity: a lambda applies a
successor, a unary operation retains the depth, and each binary or application
branch takes the original maximum with its fixed sibling's depth.
\lean{Copying/TranslatedComposition}{35}{38}{TermContext.depthAt_compose}

\subsection{Composing the unchanged translated hosts}
Choose annotated contexts $D$ and $C$, an argument allowance $d$, and arbitrary
inner and outer hole names. Translate $D$ at $d$. The universal filling theorem
provides a certified graph context $B$, an increment $n$, and a filling law
for every graph of allowance $d$. Write $q$ for its outer allowance. Translate
$C$ at $q$; the same theorem supplies a certified context $A$ and increment
$m$. The outer hole of $A$ accepts allowance $q+m$, so the actual composite
is $A\circ\operatorname{pad}_m(B)$. Its remaining hole accepts allowance
$(d+n)+m$. Both $A$ and $B$ have literally the original translated host graphs;
the construction changes neither their translation clauses nor their interfaces.

There are two simultaneous conclusions. First, apply the outer filling law to
the graph of $B$ itself. Constructor interpretation of $C$ on the translated
host of $D$ is exactly translation of $C\circ D$, by the previously proved
\texttt{graph\_compose}. Hence the composite host is recursively isomorphic
to the directly translated composed host, at the same allowance.

Second, fix any argument graph $G$. Associativity rewrites filling the
composite as filling $A$ with the filled padded $B$. The padding/filling
isomorphism identifies that inner result with padding $B$ filled by
$\operatorname{pad}_n(G)$. The outer filling law now applies. Transport the
inner filling law through the original constructor interpretation of $C$,
using its isomorphism preservation theorem. Finally, the literal
\texttt{apply\_compose} identity identifies the endpoint with interpretation
of $C\circ D$ on $G$.

The same increments, hosts and addresses work for every $G$: they are chosen
before the universal quantifier. In particular, the proof is not restricted
to language graphs or closed programs. No inequality between the two hole
names is required. The sole-hole certificates account for removal of the
outer hole before insertion of the inner context.
\lean{Copying/TranslatedComposition}{46}{78}{ContextTranslation.composed_context_filling}

\subsection{Successive filling of terms}
Supply the preceding theorem with the original graph translation of a term
$t$ at its minimal allowance. The existing exact interpretation identity
\texttt{apply\_ofTerm} and syntactic identity \texttt{fill\_compose} identify
the final graph with the original translation of $C[D[t]]$, up to recursive
isomorphism at the same outer depth. This corollary covers arbitrary finite
variable and atom scopes, including holes beneath binders, and preserves all
original maximum-depth padding choices.
\lean{Copying/TranslatedComposition}{82}{102}{ContextTranslation.composed_term_filling}

This closes the pending context-translation composition law in Appendix~B.1.
It does not supply graph-level term renaming or substitution, machine-step
invariants, the concrete copying execution or the eleven template proofs.
Proposition~12.2 remains the endpoint of those outstanding tasks.
